% --- Archon physics formalization source begin ---
% archon:physics
% archon:covers PhyXMiniProblems/problem_phyx_mini_0614.lean
% archon:source-report reports/phyx_mini/problem_phyx_mini_0614.source.json

\chapter{Physics problem phyx\_mini\_0614}
\label{ch:PhyXMiniProblems_problem_phyx_mini_0614}

\paragraph{Problem source.}
\#\# Physical scenario

An advanced civilization uses a fusion reactor to maintain a high temperature on a spherical blackbody whose radiation supplies energy for use on a spacecraft. The blackbody has radius \$r = 1.23\$ m and is surrounded by a larger sphere whose inside surface is densely lined with photocells (figure). These capture all photons that have energy within 1.00\textbackslash{}\% of a central value \$E\_0\$. The photocells are designed so that \$E\_0\$ corresponds to the peak energy radiated from the blackbody. The range \$\textbackslash{}Delta E\$ is sufficiently small that we can approximate the spectral emittance as constant over the range \$E\_0 \textbackslash{}pm \textbackslash{}Delta E\$. This device generates an emf across a motor of resistance 1.00 k\$\textbackslash{}Omega\$.The device is designed for \$T = 5000\$ K

\#\# Question

What is the current?

\#\# Answer choices

- A: A: 91.1A

- B: B: 93.3A

- C: C: 86.2A

- D: D: 88.6A

\#\# Figure caption (auxiliary; use the image as primary evidence)

The image is a diagram illustrating a blackbody and photocells setup. It features the following components:

1. **Central Sphere**: A gray-colored sphere at the center, labeled as "Blackbody, radius r." This symbolizes a blackbody with a specific radius.

2. **Radiation Arrows**: Purple wavy arrows emanate outward from the surface of the blackbody, indicating radiation being emitted in various directions.

3. **Surrounding Ring**: A ring is encircling the blackbody, representing "Photocells." This ring likely indicates a shell or arrangement of photocells designed to capture the emitted radiation.

4. **Text and Labels**: The diagram includes text labels with arrows pointing to the corresponding components. The text "Blackbody, radius r" refers to the central sphere, and "Photocells" refers to the surrounding ring.

Overall, the diagram visually represents a blackbody emitting radiation with surrounding photocells intended to detect or measure the emissions.

\#\# Dataset metadata

Quantum Phenomena; Physical Model Grounding Reasoning, Multi-Formula Reasoning

\paragraph{Recorded answer/context.}
A: A: 91.1A

\paragraph{Figure/image path.}
/root/proposal\_for\_physic/science-mango/phyx\_mini\_run/phyx\_data/test\_image/614.png

\paragraph{Formalization target.}
create a compiling Lean file with sorry bodies at `PhyXMiniProblems/problem\_phyx\_mini\_0614.lean`. The Lean declarations must preserve the physical quantities, dimensions or dimensional roles, figure labels, governing-law hypotheses, and final relation expressed by this problem.
Use Mathlib/Physlib names found through LeanExplore where available. If a physics API is missing, introduce faithful local abstractions rather than scalar placeholder aliases.

\begin{theorem}[Physics formalization target]
\label{thm:physics:phyx_mini_0614:target}
\lean{PhyXMiniProblems.ProblemPhyXMini0614.motorCurrent_matches_answer_A}
\uses{def:physics:phyx-mini-0614:phyxminiproblems-problemphyxmini0614-currentinamperes, def:physics:phyx-mini-0614:phyxminiproblems-problemphyxmini0614-blackbodyphotocellmotorsetup, def:physics:phyx-mini-0614:phyxminiproblems-problemphyxmini0614-matchesproblemstatement, def:physics:phyx-mini-0614:phyxminiproblems-problemphyxmini0614-matchesprimaryblackbodyphotocellfigure, def:physics:phyx-mini-0614:phyxminiproblems-problemphyxmini0614-usesstandardradiationconstants, def:physics:phyx-mini-0614:phyxminiproblems-problemphyxmini0614-hasphysicalblackbodyphotocellparameters, def:physics:phyx-mini-0614:phyxminiproblems-problemphyxmini0614-satisfiesblackbodyphotocellmotorlaws, def:physics:phyx-mini-0614:phyxminiproblems-problemphyxmini0614-answerchoice, def:physics:phyx-mini-0614:phyxminiproblems-problemphyxmini0614-displayedanswercurrentinamperes}
The assigned autoformalize agent should translate this physics problem into Lean declarations in the covered file, with theorem and lemma proof bodies written as `by sorry`.
\end{theorem}
\begin{proof}
This is an autoformalization task, not a proof task. Produce statements that can later be proved by the physics prover without weakening the physical model.
\end{proof}

% --- Archon declaration topology begin ---
\section{Declaration topology}
\begin{definition}[Length Quantity]
  \label{def:physics:phyx-mini-0614:phyxminiproblems-problemphyxmini0614-lengthquantity}
  \lean{PhyXMiniProblems.ProblemPhyXMini0614.LengthQuantity}
  A nonnegative physical length
\end{definition}
\begin{proof}
  This is the declared model definition of Length Quantity.
\end{proof}
\begin{definition}[Power Quantity]
  \label{def:physics:phyx-mini-0614:phyxminiproblems-problemphyxmini0614-powerquantity}
  \lean{PhyXMiniProblems.ProblemPhyXMini0614.PowerQuantity}
  Radiant or electrical power, with SI unit watt
\end{definition}
\begin{proof}
  This is the declared model definition of Power Quantity.
\end{proof}
\begin{definition}[Spectral Radiant Exitance Per Energy Quantity]
  \label{def:physics:phyx-mini-0614:phyxminiproblems-problemphyxmini0614-spectralradiantexitanceperenergyquantity}
  \lean{PhyXMiniProblems.ProblemPhyXMini0614.SpectralRadiantExitancePerEnergyQuantity}
  Spectral radiant exitance per unit photon energy, with SI unit W / (m² J) and hence physical dimension L⁻² T⁻¹
\end{definition}
\begin{proof}
  This is the declared model definition of Spectral Radiant Exitance Per Energy Quantity.
\end{proof}
\begin{definition}[Planck Constant Quantity]
  \label{def:physics:phyx-mini-0614:phyxminiproblems-problemphyxmini0614-planckconstantquantity}
  \lean{PhyXMiniProblems.ProblemPhyXMini0614.PlanckConstantQuantity}
  Planck's constant, with SI unit joule-second
\end{definition}
\begin{proof}
  This is the declared model definition of Planck Constant Quantity.
\end{proof}
\begin{definition}[Boltzmann Constant Quantity]
  \label{def:physics:phyx-mini-0614:phyxminiproblems-problemphyxmini0614-boltzmannconstantquantity}
  \lean{PhyXMiniProblems.ProblemPhyXMini0614.BoltzmannConstantQuantity}
  Boltzmann's constant, with SI unit joule per kelvin
\end{definition}
\begin{proof}
  This is the declared model definition of Boltzmann Constant Quantity.
\end{proof}
\begin{definition}[Speed Quantity]
  \label{def:physics:phyx-mini-0614:phyxminiproblems-problemphyxmini0614-speedquantity}
  \lean{PhyXMiniProblems.ProblemPhyXMini0614.SpeedQuantity}
  A physical speed. This uses a real carrier to match the actual type of Physlib's dimensionful DimSpeed.speedOfLight constant
\end{definition}
\begin{proof}
  This is the declared model definition of Speed Quantity.
\end{proof}
\begin{definition}[Electric Current Quantity]
  \label{def:physics:phyx-mini-0614:phyxminiproblems-problemphyxmini0614-electriccurrentquantity}
  \lean{PhyXMiniProblems.ProblemPhyXMini0614.ElectricCurrentQuantity}
  Electric current, with SI unit ampere and dimension charge per time
\end{definition}
\begin{proof}
  This is the declared model definition of Electric Current Quantity.
\end{proof}
\begin{definition}[Electric Potential Quantity]
  \label{def:physics:phyx-mini-0614:phyxminiproblems-problemphyxmini0614-electricpotentialquantity}
  \lean{PhyXMiniProblems.ProblemPhyXMini0614.ElectricPotentialQuantity}
  Electric potential difference or emf, with SI unit volt
\end{definition}
\begin{proof}
  This is the declared model definition of Electric Potential Quantity.
\end{proof}
\begin{definition}[Electrical Resistance Quantity]
  \label{def:physics:phyx-mini-0614:phyxminiproblems-problemphyxmini0614-electricalresistancequantity}
  \lean{PhyXMiniProblems.ProblemPhyXMini0614.ElectricalResistanceQuantity}
  Electrical resistance, with SI unit ohm
\end{definition}
\begin{proof}
  This is the declared model definition of Electrical Resistance Quantity.
\end{proof}
\begin{definition}[length In Meters]
  \label{def:physics:phyx-mini-0614:phyxminiproblems-problemphyxmini0614-lengthinmeters}
  \lean{PhyXMiniProblems.ProblemPhyXMini0614.lengthInMeters}
  \uses{def:physics:phyx-mini-0614:phyxminiproblems-problemphyxmini0614-lengthquantity}
  Read a physical length in metres
\end{definition}
\begin{proof}
  This is the declared model definition of length In Meters.
\end{proof}
\begin{definition}[energy In Joules]
  \label{def:physics:phyx-mini-0614:phyxminiproblems-problemphyxmini0614-energyinjoules}
  \lean{PhyXMiniProblems.ProblemPhyXMini0614.energyInJoules}
  Read a Physlib energy in joules
\end{definition}
\begin{proof}
  This is the declared model definition of energy In Joules.
\end{proof}
\begin{definition}[power In Watts]
  \label{def:physics:phyx-mini-0614:phyxminiproblems-problemphyxmini0614-powerinwatts}
  \lean{PhyXMiniProblems.ProblemPhyXMini0614.powerInWatts}
  \uses{def:physics:phyx-mini-0614:phyxminiproblems-problemphyxmini0614-powerquantity}
  Read a physical power in watts
\end{definition}
\begin{proof}
  This is the declared model definition of power In Watts.
\end{proof}
\begin{definition}[spectral Radiant Exitance Per Joule In SI]
  \label{def:physics:phyx-mini-0614:phyxminiproblems-problemphyxmini0614-spectralradiantexitanceperjouleinsi}
  \lean{PhyXMiniProblems.ProblemPhyXMini0614.spectralRadiantExitancePerJouleInSI}
  \uses{def:physics:phyx-mini-0614:phyxminiproblems-problemphyxmini0614-spectralradiantexitanceperenergyquantity}
  Read spectral radiant exitance in W / (m² J)
\end{definition}
\begin{proof}
  This is the declared model definition of spectral Radiant Exitance Per Joule In SI.
\end{proof}
\begin{definition}[planck Constant In Joule Seconds]
  \label{def:physics:phyx-mini-0614:phyxminiproblems-problemphyxmini0614-planckconstantinjouleseconds}
  \lean{PhyXMiniProblems.ProblemPhyXMini0614.planckConstantInJouleSeconds}
  \uses{def:physics:phyx-mini-0614:phyxminiproblems-problemphyxmini0614-planckconstantquantity}
  Read Planck's constant in joule-seconds
\end{definition}
\begin{proof}
  This is the declared model definition of planck Constant In Joule Seconds.
\end{proof}
\begin{definition}[boltzmann Constant In Joules Per Kelvin]
  \label{def:physics:phyx-mini-0614:phyxminiproblems-problemphyxmini0614-boltzmannconstantinjoulesperkelvin}
  \lean{PhyXMiniProblems.ProblemPhyXMini0614.boltzmannConstantInJoulesPerKelvin}
  \uses{def:physics:phyx-mini-0614:phyxminiproblems-problemphyxmini0614-boltzmannconstantquantity}
  Read Boltzmann's constant in joules per kelvin
\end{definition}
\begin{proof}
  This is the declared model definition of boltzmann Constant In Joules Per Kelvin.
\end{proof}
\begin{definition}[speed In Meters Per Second]
  \label{def:physics:phyx-mini-0614:phyxminiproblems-problemphyxmini0614-speedinmeterspersecond}
  \lean{PhyXMiniProblems.ProblemPhyXMini0614.speedInMetersPerSecond}
  \uses{def:physics:phyx-mini-0614:phyxminiproblems-problemphyxmini0614-speedquantity}
  Read a dimensionful speed in metres per second
\end{definition}
\begin{proof}
  This is the declared model definition of speed In Meters Per Second.
\end{proof}
\begin{definition}[current In Amperes]
  \label{def:physics:phyx-mini-0614:phyxminiproblems-problemphyxmini0614-currentinamperes}
  \lean{PhyXMiniProblems.ProblemPhyXMini0614.currentInAmperes}
  \uses{def:physics:phyx-mini-0614:phyxminiproblems-problemphyxmini0614-electriccurrentquantity}
  Read electric current in amperes
\end{definition}
\begin{proof}
  This is the declared model definition of current In Amperes.
\end{proof}
\begin{definition}[potential In Volts]
  \label{def:physics:phyx-mini-0614:phyxminiproblems-problemphyxmini0614-potentialinvolts}
  \lean{PhyXMiniProblems.ProblemPhyXMini0614.potentialInVolts}
  \uses{def:physics:phyx-mini-0614:phyxminiproblems-problemphyxmini0614-electricpotentialquantity}
  Read emf or potential difference in volts
\end{definition}
\begin{proof}
  This is the declared model definition of potential In Volts.
\end{proof}
\begin{definition}[resistance In Ohms]
  \label{def:physics:phyx-mini-0614:phyxminiproblems-problemphyxmini0614-resistanceinohms}
  \lean{PhyXMiniProblems.ProblemPhyXMini0614.resistanceInOhms}
  \uses{def:physics:phyx-mini-0614:phyxminiproblems-problemphyxmini0614-electricalresistancequantity}
  Read electrical resistance in ohms
\end{definition}
\begin{proof}
  This is the declared model definition of resistance In Ohms.
\end{proof}
\begin{definition}[temperature In Kelvins]
  \label{def:physics:phyx-mini-0614:phyxminiproblems-problemphyxmini0614-temperatureinkelvins}
  \lean{PhyXMiniProblems.ProblemPhyXMini0614.temperatureInKelvins}
  Read a Physlib absolute temperature in kelvins. The storage unit is explicit because Temperature stores a nonnegative magnitude in an arbitrary zero-preserving temperature unit
\end{definition}
\begin{proof}
  This is the declared model definition of temperature In Kelvins.
\end{proof}
\begin{definition}[Thermal Emitter Model]
  \label{def:physics:phyx-mini-0614:phyxminiproblems-problemphyxmini0614-thermalemittermodel}
  \lean{PhyXMiniProblems.ProblemPhyXMini0614.ThermalEmitterModel}
  The radiation model assigned to the central hot sphere
\end{definition}
\begin{proof}
  This is the declared model definition of Thermal Emitter Model.
\end{proof}
\begin{definition}[Figure Component Label]
  \label{def:physics:phyx-mini-0614:phyxminiproblems-problemphyxmini0614-figurecomponentlabel}
  \lean{PhyXMiniProblems.ProblemPhyXMini0614.FigureComponentLabel}
  Component names printed in the supplied figure
\end{definition}
\begin{proof}
  This is the declared model definition of Figure Component Label.
\end{proof}
\begin{definition}[Figure Radius Symbol]
  \label{def:physics:phyx-mini-0614:phyxminiproblems-problemphyxmini0614-figureradiussymbol}
  \lean{PhyXMiniProblems.ProblemPhyXMini0614.FigureRadiusSymbol}
  The radius symbol printed beside the central blackbody
\end{definition}
\begin{proof}
  This is the declared model definition of Figure Radius Symbol.
\end{proof}
\begin{definition}[Blackbody Photocell Figure]
  \label{def:physics:phyx-mini-0614:phyxminiproblems-problemphyxmini0614-blackbodyphotocellfigure}
  \lean{PhyXMiniProblems.ProblemPhyXMini0614.BlackbodyPhotocellFigure}
  \uses{def:physics:phyx-mini-0614:phyxminiproblems-problemphyxmini0614-figurecomponentlabel, def:physics:phyx-mini-0614:phyxminiproblems-problemphyxmini0614-figureradiussymbol}
  Typed information carried by image 614. The quantitative radius is supplied by the prose; the bitmap contributes the component placement, labels, radius symbol, enclosing photocell lining, and outward radiation arrows
\end{definition}
\begin{proof}
  This is the declared model definition of Blackbody Photocell Figure.
\end{proof}
\begin{definition}[Spherical Blackbody]
  \label{def:physics:phyx-mini-0614:phyxminiproblems-problemphyxmini0614-sphericalblackbody}
  \lean{PhyXMiniProblems.ProblemPhyXMini0614.SphericalBlackbody}
  \uses{def:physics:phyx-mini-0614:phyxminiproblems-problemphyxmini0614-lengthquantity, def:physics:phyx-mini-0614:phyxminiproblems-problemphyxmini0614-spectralradiantexitanceperenergyquantity, def:physics:phyx-mini-0614:phyxminiproblems-problemphyxmini0614-thermalemittermodel}
  The spherical source. Its spectral radiant exitance is an independent physical observable indexed by a genuine physical photon energy; Planck's law is imposed separately below
\end{definition}
\begin{proof}
  This is the declared model definition of Spherical Blackbody.
\end{proof}
\begin{definition}[Photocell Shell]
  \label{def:physics:phyx-mini-0614:phyxminiproblems-problemphyxmini0614-photocellshell}
  \lean{PhyXMiniProblems.ProblemPhyXMini0614.PhotocellShell}
  \uses{def:physics:phyx-mini-0614:phyxminiproblems-problemphyxmini0614-powerquantity}
  The enclosing photocell shell. capturesPhotonAtEnergy records its response window, while captured power remains independent until the narrow-band law is assumed
\end{definition}
\begin{proof}
  This is the declared model definition of Photocell Shell.
\end{proof}
\begin{definition}[Resistive Motor]
  \label{def:physics:phyx-mini-0614:phyxminiproblems-problemphyxmini0614-resistivemotor}
  \lean{PhyXMiniProblems.ProblemPhyXMini0614.ResistiveMotor}
  \uses{def:physics:phyx-mini-0614:phyxminiproblems-problemphyxmini0614-powerquantity, def:physics:phyx-mini-0614:phyxminiproblems-problemphyxmini0614-electriccurrentquantity, def:physics:phyx-mini-0614:phyxminiproblems-problemphyxmini0614-electricpotentialquantity, def:physics:phyx-mini-0614:phyxminiproblems-problemphyxmini0614-electricalresistancequantity}
  The motor's independent circuit observables
\end{definition}
\begin{proof}
  This is the declared model definition of Resistive Motor.
\end{proof}
\begin{definition}[Blackbody Radiation Constants]
  \label{def:physics:phyx-mini-0614:phyxminiproblems-problemphyxmini0614-blackbodyradiationconstants}
  \lean{PhyXMiniProblems.ProblemPhyXMini0614.BlackbodyRadiationConstants}
  \uses{def:physics:phyx-mini-0614:phyxminiproblems-problemphyxmini0614-planckconstantquantity, def:physics:phyx-mini-0614:phyxminiproblems-problemphyxmini0614-boltzmannconstantquantity, def:physics:phyx-mini-0614:phyxminiproblems-problemphyxmini0614-speedquantity}
  Dimensionful constants occurring in Planck's law
\end{definition}
\begin{proof}
  This is the declared model definition of Blackbody Radiation Constants.
\end{proof}
\begin{definition}[Blackbody Photocell Motor Setup]
  \label{def:physics:phyx-mini-0614:phyxminiproblems-problemphyxmini0614-blackbodyphotocellmotorsetup}
  \lean{PhyXMiniProblems.ProblemPhyXMini0614.BlackbodyPhotocellMotorSetup}
  \uses{def:physics:phyx-mini-0614:phyxminiproblems-problemphyxmini0614-blackbodyphotocellfigure, def:physics:phyx-mini-0614:phyxminiproblems-problemphyxmini0614-sphericalblackbody, def:physics:phyx-mini-0614:phyxminiproblems-problemphyxmini0614-photocellshell, def:physics:phyx-mini-0614:phyxminiproblems-problemphyxmini0614-resistivemotor, def:physics:phyx-mini-0614:phyxminiproblems-problemphyxmini0614-blackbodyradiationconstants}
  Complete physical and diagrammatic setup for the problem
\end{definition}
\begin{proof}
  This is the declared model definition of Blackbody Photocell Motor Setup.
\end{proof}
\begin{definition}[Matches Problem Statement]
  \label{def:physics:phyx-mini-0614:phyxminiproblems-problemphyxmini0614-matchesproblemstatement}
  \lean{PhyXMiniProblems.ProblemPhyXMini0614.MatchesProblemStatement}
  \uses{def:physics:phyx-mini-0614:phyxminiproblems-problemphyxmini0614-lengthinmeters, def:physics:phyx-mini-0614:phyxminiproblems-problemphyxmini0614-energyinjoules, def:physics:phyx-mini-0614:phyxminiproblems-problemphyxmini0614-spectralradiantexitanceperjouleinsi, def:physics:phyx-mini-0614:phyxminiproblems-problemphyxmini0614-resistanceinohms, def:physics:phyx-mini-0614:phyxminiproblems-problemphyxmini0614-temperatureinkelvins, def:physics:phyx-mini-0614:phyxminiproblems-problemphyxmini0614-blackbodyphotocellmotorsetup}
  Numerical data and qualitative facts supplied by the prose. In particular, the central accepted energy maximizes the spectrum, but neither its numerical value nor the requested motor current is assumed
\end{definition}
\begin{proof}
  This is the declared model definition of Matches Problem Statement.
\end{proof}
\begin{definition}[Matches Primary Blackbody Photocell Figure]
  \label{def:physics:phyx-mini-0614:phyxminiproblems-problemphyxmini0614-matchesprimaryblackbodyphotocellfigure}
  \lean{PhyXMiniProblems.ProblemPhyXMini0614.MatchesPrimaryBlackbodyPhotocellFigure}
  \uses{def:physics:phyx-mini-0614:phyxminiproblems-problemphyxmini0614-blackbodyphotocellmotorsetup}
  Primary-image evidence: the blackbody is central, the photocells surround it, both names and the radius symbol r are shown, and eight wavy radiation arrows point from the source toward the shell
\end{definition}
\begin{proof}
  This is the declared model definition of Matches Primary Blackbody Photocell Figure.
\end{proof}
\begin{definition}[Uses Standard Radiation Constants]
  \label{def:physics:phyx-mini-0614:phyxminiproblems-problemphyxmini0614-usesstandardradiationconstants}
  \lean{PhyXMiniProblems.ProblemPhyXMini0614.UsesStandardRadiationConstants}
  \uses{def:physics:phyx-mini-0614:phyxminiproblems-problemphyxmini0614-planckconstantinjouleseconds, def:physics:phyx-mini-0614:phyxminiproblems-problemphyxmini0614-boltzmannconstantinjoulesperkelvin, def:physics:phyx-mini-0614:phyxminiproblems-problemphyxmini0614-blackbodyphotocellmotorsetup}
  Standard SI calibrations. Physlib supplies the reduced Planck constant and the dimensionful speed of light; the full Planck constant is h = 2πℏ
\end{definition}
\begin{proof}
  This is the declared model definition of Uses Standard Radiation Constants.
\end{proof}
\begin{definition}[Has Physical Blackbody Photocell Parameters]
  \label{def:physics:phyx-mini-0614:phyxminiproblems-problemphyxmini0614-hasphysicalblackbodyphotocellparameters}
  \lean{PhyXMiniProblems.ProblemPhyXMini0614.HasPhysicalBlackbodyPhotocellParameters}
  \uses{def:physics:phyx-mini-0614:phyxminiproblems-problemphyxmini0614-lengthinmeters, def:physics:phyx-mini-0614:phyxminiproblems-problemphyxmini0614-energyinjoules, def:physics:phyx-mini-0614:phyxminiproblems-problemphyxmini0614-planckconstantinjouleseconds, def:physics:phyx-mini-0614:phyxminiproblems-problemphyxmini0614-boltzmannconstantinjoulesperkelvin, def:physics:phyx-mini-0614:phyxminiproblems-problemphyxmini0614-speedinmeterspersecond, def:physics:phyx-mini-0614:phyxminiproblems-problemphyxmini0614-resistanceinohms, def:physics:phyx-mini-0614:phyxminiproblems-problemphyxmini0614-temperatureinkelvins, def:physics:phyx-mini-0614:phyxminiproblems-problemphyxmini0614-blackbodyphotocellmotorsetup}
  Positivity and nondegeneracy conditions for the physical apparatus
\end{definition}
\begin{proof}
  This is the declared model definition of Has Physical Blackbody Photocell Parameters.
\end{proof}
\begin{definition}[Satisfies Blackbody Photocell Motor Laws]
  \label{def:physics:phyx-mini-0614:phyxminiproblems-problemphyxmini0614-satisfiesblackbodyphotocellmotorlaws}
  \lean{PhyXMiniProblems.ProblemPhyXMini0614.SatisfiesBlackbodyPhotocellMotorLaws}
  \uses{def:physics:phyx-mini-0614:phyxminiproblems-problemphyxmini0614-lengthinmeters, def:physics:phyx-mini-0614:phyxminiproblems-problemphyxmini0614-energyinjoules, def:physics:phyx-mini-0614:phyxminiproblems-problemphyxmini0614-powerinwatts, def:physics:phyx-mini-0614:phyxminiproblems-problemphyxmini0614-spectralradiantexitanceperjouleinsi, def:physics:phyx-mini-0614:phyxminiproblems-problemphyxmini0614-planckconstantinjouleseconds, def:physics:phyx-mini-0614:phyxminiproblems-problemphyxmini0614-boltzmannconstantinjoulesperkelvin, def:physics:phyx-mini-0614:phyxminiproblems-problemphyxmini0614-speedinmeterspersecond, def:physics:phyx-mini-0614:phyxminiproblems-problemphyxmini0614-currentinamperes, def:physics:phyx-mini-0614:phyxminiproblems-problemphyxmini0614-potentialinvolts, def:physics:phyx-mini-0614:phyxminiproblems-problemphyxmini0614-resistanceinohms, def:physics:phyx-mini-0614:phyxminiproblems-problemphyxmini0614-temperatureinkelvins, def:physics:phyx-mini-0614:phyxminiproblems-problemphyxmini0614-blackbodyphotocellmotorsetup}
  Governing physical laws. The first field is Planck's blackbody spectral radiant exitance per unit photon energy. The second implements exactly the narrow-window approximation stated in the problem. The final three fields model ideal energy conversion and the ordinary resistive circuit. None of these laws mentions 91.1 A or any answer label
\end{definition}
\begin{proof}
  This is the declared model definition of Satisfies Blackbody Photocell Motor Laws.
\end{proof}
\begin{lemma}[central Accepted Energy satisfies wien Peak Equation]
  \label{lem:physics:phyx-mini-0614:phyxminiproblems-problemphyxmini0614-centralacceptedenergy-satisfies-wienpeakequation}
  \lean{PhyXMiniProblems.ProblemPhyXMini0614.centralAcceptedEnergy_satisfies_wienPeakEquation}
  \uses{def:physics:phyx-mini-0614:phyxminiproblems-problemphyxmini0614-energyinjoules, def:physics:phyx-mini-0614:phyxminiproblems-problemphyxmini0614-boltzmannconstantinjoulesperkelvin, def:physics:phyx-mini-0614:phyxminiproblems-problemphyxmini0614-temperatureinkelvins, def:physics:phyx-mini-0614:phyxminiproblems-problemphyxmini0614-blackbodyphotocellmotorsetup, def:physics:phyx-mini-0614:phyxminiproblems-problemphyxmini0614-matchesproblemstatement, def:physics:phyx-mini-0614:phyxminiproblems-problemphyxmini0614-hasphysicalblackbodyphotocellparameters, def:physics:phyx-mini-0614:phyxminiproblems-problemphyxmini0614-satisfiesblackbodyphotocellmotorlaws}
  The energy-spectrum maximum satisfies the dimensionless Wien peak equation 3 (1 - exp (-x₀)) = x₀, where x₀ = E₀/(k\_B T). This is derived from the given peak condition and Planck's law; no numerical root is assumed
\end{lemma}
\begin{proof}
  The conclusion follows from the governing relations and figure readouts named in the statement of central Accepted Energy satisfies wien Peak Equation.
\end{proof}
\begin{lemma}[motor Current squared times resistance eq narrow Band Power]
  \label{lem:physics:phyx-mini-0614:phyxminiproblems-problemphyxmini0614-motorcurrent-squared-times-resistance-eq-narrowbandpower}
  \lean{PhyXMiniProblems.ProblemPhyXMini0614.motorCurrent_squared_times_resistance_eq_narrowBandPower}
  \uses{def:physics:phyx-mini-0614:phyxminiproblems-problemphyxmini0614-lengthinmeters, def:physics:phyx-mini-0614:phyxminiproblems-problemphyxmini0614-energyinjoules, def:physics:phyx-mini-0614:phyxminiproblems-problemphyxmini0614-spectralradiantexitanceperjouleinsi, def:physics:phyx-mini-0614:phyxminiproblems-problemphyxmini0614-currentinamperes, def:physics:phyx-mini-0614:phyxminiproblems-problemphyxmini0614-resistanceinohms, def:physics:phyx-mini-0614:phyxminiproblems-problemphyxmini0614-blackbodyphotocellmotorsetup, def:physics:phyx-mini-0614:phyxminiproblems-problemphyxmini0614-satisfiesblackbodyphotocellmotorlaws}
  Combining the narrow-band radiation approximation with ideal conversion, P = V I, and Ohm's law yields the current-squared power balance. This is a derived relation, not a premise of the final numerical result
\end{lemma}
\begin{proof}
  The conclusion follows from the governing relations and figure readouts named in the statement of motor Current squared times resistance eq narrow Band Power.
\end{proof}
\begin{definition}[Answer Choice]
  \label{def:physics:phyx-mini-0614:phyxminiproblems-problemphyxmini0614-answerchoice}
  \lean{PhyXMiniProblems.ProblemPhyXMini0614.AnswerChoice}
  Labels of the four displayed answers
\end{definition}
\begin{proof}
  This is the declared model definition of Answer Choice.
\end{proof}
\begin{definition}[displayed Answer Current In Amperes]
  \label{def:physics:phyx-mini-0614:phyxminiproblems-problemphyxmini0614-displayedanswercurrentinamperes}
  \lean{PhyXMiniProblems.ProblemPhyXMini0614.displayedAnswerCurrentInAmperes}
  \uses{def:physics:phyx-mini-0614:phyxminiproblems-problemphyxmini0614-answerchoice}
  Current in amperes printed beside each answer label
\end{definition}
\begin{proof}
  This is the declared model definition of displayed Answer Current In Amperes.
\end{proof}
\begin{definition}[recorded Dataset Answer]
  \label{def:physics:phyx-mini-0614:phyxminiproblems-problemphyxmini0614-recordeddatasetanswer}
  \lean{PhyXMiniProblems.ProblemPhyXMini0614.recordedDatasetAnswer}
  \uses{def:physics:phyx-mini-0614:phyxminiproblems-problemphyxmini0614-answerchoice}
  The dataset records choice A; this metadata is not accepted as a premise
\end{definition}
\begin{proof}
  This is the declared model definition of recorded Dataset Answer.
\end{proof}
% --- Archon declaration topology end ---
% --- Archon physics formalization source end ---
